\hypertarget{chapter-32-c23-core-language}{%
\section{CHAPTER 32: C23 CORE
LANGUAGE}\label{chapter-32-c23-core-language}}

\hypertarget{c23-core-language-upgrades}{%
\section{C++23 CORE LANGUAGE
UPGRADES}\label{c23-core-language-upgrades}}

C++23 is the ``Ergonomics'' release. It polishes the massive changes
introduced in C++20, removing boilerplate and completing the modern C++
paradigm.

\hypertarget{deducing-this-explicit-object-parameters}{%
\subsubsection{\texorpdfstring{1. Deducing \texttt{this} (Explicit
Object
Parameters)}{1. Deducing this (Explicit Object Parameters)}}\label{deducing-this-explicit-object-parameters}}

One of the most revolutionary changes for class design. It allows a
member function to explicitly declare the object it is called on as its
first parameter. * \textbf{Recursive Lambdas}: A lambda can now easily
call itself without \passthrough{\lstinline!std::function!} overhead.
\passthrough{\lstinline!cpp     auto fib = [](this auto self, int n) -> int \{         return n <= 1 ? n : self(n - 1) + self(n - 2);     \};!}
* \textbf{Simplifying CRTP \& Forwarding}: You no longer need 4
overloads (\passthrough{\lstinline!\&!},
\passthrough{\lstinline!const\&!}, \passthrough{\lstinline!\&\&!},
\passthrough{\lstinline!const\&\&!}) or complex CRTP inheritance to
perfectly forward a member.
\passthrough{\lstinline!cpp     template<typename Self>     auto\&\& get\_data(this Self\&\& self) \{         // Automatically preserves const/ref qualifiers of the object         return std::forward<Self>(self).data;      \}!}

\hypertarget{syntactic-ergonomics}{%
\subsubsection{2. Syntactic Ergonomics}\label{syntactic-ergonomics}}

\begin{itemize}
\tightlist
\item
  \textbf{Multidimensional Subscript Operator}: You can now pass
  multiple arguments to \passthrough{\lstinline!operator[]!}, which is
  crucial for linear algebra (pairs perfectly with
  \passthrough{\lstinline!std::mdspan!}).
  \passthrough{\lstinline!cpp     struct Matrix \{         double\& operator[](size\_t row, size\_t col) \{ return data[row * cols + col]; \}     \};     matrix[1, 2] = 42.0;!}
\item
  \textbf{\passthrough{\lstinline!if consteval!}}: A cleaner,
  standardized way to execute different code depending on whether the
  function is evaluated at compile time or run time.
  \passthrough{\lstinline!cpp     constexpr double power(double d, int p) \{         if consteval \{ return compile\_time\_pow(d, p); \}         else \{ return std::pow(d, p); \} // Run time     \}!}
\item
  \textbf{\passthrough{\lstinline!auto(x)!} and
  \passthrough{\lstinline!auto\{x\}!} (Decay Copy)}: Explicitly requests
  a PR-value copy of a variable, forcing decay (useful in generic macros
  and templates).
  \passthrough{\lstinline!cpp     void f(auto\& x) \{         auto copy = auto(x); // Explicit copy, stripping references     \}!}
\item
  \textbf{Static \passthrough{\lstinline!operator()!}}: Lambdas that
  don't capture anything can now have a \passthrough{\lstinline!static!}
  call operator, allowing them to be passed as traditional C function
  pointers seamlessly.
  \passthrough{\lstinline!cpp     auto f = [] static (int x) \{ return x * 2; \};!}
\item
  \textbf{Size\_t Literal Suffixes}: Use \passthrough{\lstinline!z!} for
  signed \passthrough{\lstinline!ssize\_t!} and
  \passthrough{\lstinline!uz!} for unsigned
  \passthrough{\lstinline!size\_t!}.
  \passthrough{\lstinline!cpp     for (auto i = 0uz; i < vec.size(); ++i) \{\}!}
\item
  \textbf{Labels at the end of compound statements}: You no longer need
  to put a dummy statement after a label at the end of a block.
\end{itemize}
